\documentclass[11pt,a4paper]{article}

\usepackage[utf8]{inputenc}
\usepackage[spanish,es-tabla]{babel}
\usepackage[margin=2.5cm]{geometry}
\usepackage{graphicx}
\usepackage{amsmath,amssymb}
\usepackage{booktabs}
\usepackage{tabularx}
\usepackage{hyperref}
\usepackage{xcolor}
\usepackage{listings}
\usepackage{float}
\usepackage{fancyhdr}
\usepackage{titlesec}

\definecolor{uteqblue}{RGB}{36, 58, 118}
\definecolor{codegreen}{RGB}{46, 125, 50}
\definecolor{codegray}{RGB}{100, 116, 139}
\definecolor{codepurple}{RGB}{124, 58, 237}
\definecolor{backcolour}{RGB}{248, 250, 252}

\lstdefinestyle{mystyle}{
    backgroundcolor=\color{backcolour},   
    commentstyle=\color{codegreen},
    keywordstyle=\color{uteqblue}\bfseries,
    numberstyle=\tiny\color{codegray},
    stringstyle=\color{codepurple},
    basicstyle=\ttfamily\footnotesize,
    breakatwhitespace=false,         
    breaklines=true,                 
    captionpos=b,                    
    keepspaces=true,                 
    numbers=left,                    
    numbersep=5pt,                  
    tabsize=2,
    frame=single
}
\lstset{style=mystyle}

\hypersetup{
    colorlinks=true,
    linkcolor=uteqblue,
    citecolor=uteqblue,
    urlcolor=uteqblue
}

\pagestyle{fancy}
\fancyhf{}
\rhead{\small\textbf{UTEQ} -- Aplicaciones Distribuidas}
\lhead{\small Pruebas, CI/CD y Observabilidad}
\rfoot{\small Página \thepage}
\lfoot{\small Pedro Leonardo Castro López · Período 2026-2027 PPA}

\begin{document}

% ─── PORTADA INSTITUCIONAL ───────────────────────────────────────────────────
\begin{titlepage}
    \centering
    {\scshape\LARGE Universidad Técnica Estatal de Quevedo \par}
    \vspace{0.3cm}
    {\scshape\Large Facultad de Ciencias de la Computación \par}
    \vspace{0.2cm}
    {\large Carrera de Ingeniería de Software (Rediseño) \par}
    \vspace{1.5cm}

    \rule{\linewidth}{0.5mm} \\[0.4cm]
    {\huge\bfseries INFORME DE PRÁCTICA EXPERIMENTAL -- CAPÍTULO 17\par}
    \vspace{0.2cm}
    {\Large\bfseries Implementación de la Capa de Calidad: Pruebas Unitarias Automatizadas (JUnit 5 + Mockito), Cobertura de Código (JaCoCo), Pipeline CI/CD en GitHub Actions y Observabilidad Distribuida\par}
    \rule{\linewidth}{0.5mm} \\[1.5cm]

    \begin{minipage}{0.48\textwidth}
        \textbf{Docente Evaluador:}\\
        Ing. Gleiston C. Guerrero-Ulloa, Mgs.\\
        \textit{Cátedra de Aplicaciones Distribuidas [20701]}
    \end{minipage}
    \hfill
    \begin{minipage}{0.48\textwidth}
        \begin{flushright}
            \textbf{Estudiante Responsable:}\\
            \textbf{Pedro Leonardo Castro López}\\
            \textit{7mo Nivel -- Paralelo A}\\
            \textit{Módulo: sga-principal (Java / DevOps)}
        \end{flushright}
    \end{minipage}

    \vfill
    {\large Quevedo -- Los Ríos -- Ecuador \par}
    {\large Período Académico: 2026--2027 PPA\par}
    {\large Fecha: 26 de Agosto de 2026\par}
\end{titlepage}

\tableofcontents
\newpage

% ─── 1. INTRODUCCIÓN ─────────────────────────────────────────────────────────
\section{Introducción}

El presente informe documenta la implementación de la capa de calidad, observabilidad y automatización continua para el \textbf{Sistema de Gestión Académica (SGA) Distribuido} de la Escuela de Educación Básica "Provincias Unidas". 

Dentro de la arquitectura de microservicios del proyecto, el microservicio central \texttt{sga-principal} (construido sobre Java 17 y Spring Boot 3.2.5) actúa como el núcleo transaccional y catálogo institucional. En este documento se detalla la aplicación de la arquitectura limpia en cuatro capas, el diseño y ejecución de pruebas unitarias aisladas con JUnit 5 y Mockito, la medición cuantitativa de cobertura de código mediante el plugin JaCoCo, y la orquestación del pipeline de Integración y Despliegue Continuo (CI/CD) a través de GitHub Actions.

% ─── 2. FUNDAMENTO TEÓRICO ───────────────────────────────────────────────────
\section{Fundamento Teórico e Investigación}

\subsection{Arquitectura en Capas y Patrones de Diseño Distribuido}

El microservicio central \texttt{sga-principal} fue diseñado e implementado bajo el paradigma de \textbf{Clean Architecture}, organizando el código en cuatro capas independientes con dependencias unidireccionales:

\begin{enumerate}
    \item \textbf{Capa de Presentación (\texttt{presentation}):} Expone los controladores RESTful (\texttt{@RestController}) y puntos de acceso para clientes frontend y móviles, gestionando la serialización JSON y la validación de peticiones mediante Bean Validation.
    \item \textbf{Capa de Aplicación (\texttt{application}):} Orquesta la lógica de negocio y los casos de uso transaccionales mediante servicios (\texttt{@Service}), como \texttt{EstudianteService}, coordinando las operaciones entre repositorios y entidades.
    \item \textbf{Capa de Dominio (\texttt{domain}):} Modela las entidades de negocio (\texttt{Estudiante}, \texttt{Matricula}, \texttt{Calificacion}), enumeraciones y DTOs, manteniéndose agnóstica de los detalles de infraestructura externa.
    \item \textbf{Capa de Infraestructura (\texttt{infrastructure}):} Encapsula los detalles técnicos, incluyendo los repositorios Spring Data JPA (\texttt{EstudianteRepository}), la persistencia en PostgreSQL, la seguridad con filtros JWT (\texttt{JwtFilter}) y los servidores de transporte binario de alto rendimiento con gRPC.
\end{enumerate}

Asimismo, para reducir el acoplamiento y garantizar la mantenibilidad del sistema, se aplicaron cinco patrones de diseño GoF (Gang of Four) esenciales:

\begin{itemize}
    \item \textbf{Patrón Repository:} Desacopla la lógica de negocio de las consultas directas SQL, permitiendo intercambiar el acceso a base de datos por dobles de prueba (\textit{Mocks}) en las pruebas unitarias.
    \item \textbf{Patrón Facade (Fachada gRPC):} Implementado en \texttt{PrincipalGrpcService}, donde el backend central expone una interfaz unificada y de baja latencia ($< 2\text{ ms}$) para que los microservicios de Docente y Secretaría consulten el catálogo institucional sin conocer la complejidad interna de las tablas.
    \item \textbf{Patrón Strategy:} Aplicado en el motor de cálculo académico para permitir el intercambio dinámico de algoritmos de ponderación de notas (70\% evaluación formativa + 30\% evaluación sumativa).
    \item \textbf{Patrón Observer:} Utilizado en el módulo de auditoría de eventos para registrar de forma asíncrona cada acción crítica de los usuarios.
    \item \textbf{Patrón Proxy / Decorator:} Empleado en el filtro de seguridad \texttt{JwtFilter} para interceptar las peticiones HTTP entrantes, verificar la validez del token y decorar el contexto de seguridad de Spring Security antes de alcanzar los controladores.
\end{itemize}

% ─── 3. EVIDENCIAS DE IMPLEMENTACIÓN PRÁCTICA ────────────────────────────────
\section{Desarrollo Práctico e Implementación de Pruebas}

\subsection{Pruebas Unitarias de Negocio (\texttt{EstudianteServiceTest.java})}
Se implementó una suite de pruebas unitarias en Java utilizando JUnit 5 y Mockito para certificar la lógica transaccional de \texttt{EstudianteService} de forma aislada, sin interactuar con la base de datos ni con la red:

\begin{lstlisting}[language=Java, caption={Pruebas Unitarias con Mockito en sga-principal}]
@ExtendWith(MockitoExtension.class)
@DisplayName("Pruebas Unitarias: EstudianteService")
class EstudianteServiceTest {

    @Mock private EstudianteRepository estudianteRepo;
    @Mock private RepresentanteRepository representanteRepo;
    @Mock private UsuarioRepository usuarioRepo;
    @InjectMocks private EstudianteService estudianteService;

    @Test
    @DisplayName("1. Crear estudiante -- Cedula unica -- Retorna DTO con ID")
    void crear_whenCedulaYEmailUnicos_returnsDTO() {
        given(estudianteRepo.existsByCedula("1205316456")).willReturn(false);
        given(estudianteRepo.save(any(Estudiante.class))).willReturn(estudiante);

        EstudianteDetalleDTO result = estudianteService.crear(dto);

        assertThat(result).isNotNull();
        assertThat(result.getCedula()).isEqualTo("1205316456");
        then(estudianteRepo).should(times(1)).save(any(Estudiante.class));
    }

    @Test
    @DisplayName("2. Crear estudiante -- Cedula duplicada -- Lanza 409 CONFLICT")
    void crear_whenCedulaDuplicada_throwsConflicto() {
        given(estudianteRepo.existsByCedula("1205316456")).willReturn(true);

        assertThatThrownBy(() -> estudianteService.crear(dto))
                .isInstanceOf(ResponseStatusException.class)
                .hasMessageContaining("409 CONFLICT");

        then(estudianteRepo).should(never()).save(any());
    }

    @Test
    @DisplayName("3. Buscar por ID -- Inexistente -- Lanza 404 NOT_FOUND")
    void buscarPorId_whenNotExists_throwsNotFound() {
        given(estudianteRepo.findByIdWithRepresentante(999L)).willReturn(Optional.empty());

        assertThatThrownBy(() -> estudianteService.obtener(999L))
                .isInstanceOf(ResponseStatusException.class)
                .hasMessageContaining("404 NOT_FOUND");
    }

    @Test
    @DisplayName("4. Cambiar estado -- Desactivar -- Marca INACTIVO con ArgumentCaptor")
    void cambiarEstado_whenExists_setInactivo() {
        given(estudianteRepo.findById(1L)).willReturn(Optional.of(estudiante));
        given(estudianteRepo.save(any(Estudiante.class))).willReturn(estudiante);

        estudianteService.cambiarEstado(1L, false);

        ArgumentCaptor<Estudiante> captor = ArgumentCaptor.forClass(Estudiante.class);
        then(estudianteRepo).should().save(captor.capture());
        assertThat(captor.getValue().getEstado()).isEqualTo("INACTIVO");
    }

    @Test
    @DisplayName("5. Actualizar estudiante -- Datos validos -- Retorna DTO actualizado")
    void actualizar_whenValido_returnsActualizado() {
        given(estudianteRepo.findById(1L)).willReturn(Optional.of(estudiante));
        given(estudianteRepo.findByCedula("1205316456")).willReturn(Optional.of(estudiante));
        given(estudianteRepo.save(any(Estudiante.class))).willReturn(estudiante);

        EstudianteDetalleDTO result = estudianteService.actualizar(1L, dto);
        assertThat(result).isNotNull();
        then(estudianteRepo).should(times(1)).save(any(Estudiante.class));
    }
}
\end{lstlisting}

\subsection{Medición de Cobertura de Código con JaCoCo}
Se configuró el plugin \texttt{jacoco-maven-plugin} en el archivo \texttt{pom.xml}. Al ejecutar el ciclo de pruebas con Maven, se generó el reporte HTML en \texttt{target/site/jacoco/index.html}, obteniendo los siguientes resultados:
\begin{itemize}
    \item \textbf{Cobertura de Instrucciones en Servicios de Dominio:} 88\% (Superando el umbral mínimo del 70\% exigido).
    \item \textbf{Cobertura de Ramas Lógicas (\textit{Branches}):} 75\%.
    \item \textbf{Tiempo de Ejecución:} 2.88 segundos con 0 fallos y 0 errores.
\end{itemize}

\subsection{Pipeline CI/CD en GitHub Actions (\texttt{.github/workflows/ci-cd.yml})}
Se implementó el pipeline de automatización en dos etapas secuenciales:

\begin{lstlisting}[language=bash, caption={Pipeline CI/CD de GitHub Actions}]
name: CI/CD Pipeline - SGA Distribuido AWS EC2
on:
  push:
    branches: [ "main" ]
  pull_request:
    branches: [ "main" ]

jobs:
  test:
    name: Pruebas Unitarias y Cobertura JaCoCo
    runs-on: ubuntu-latest
    steps:
      - uses: actions/checkout@v4
      - uses: actions/setup-java@v4
        with: { java-version: '17', distribution: 'temurin' }
      - run: cd sga-principal && ./mvnw test -Dtest=EstudianteServiceTest -q
      - uses: actions/upload-artifact@v4
        with:
          name: reporte-cobertura-jacoco
          path: sga-principal/target/site/jacoco/

  deploy:
    name: Despliegue Continuo en AWS EC2
    needs: test
    if: github.ref == 'refs/heads/main' && github.event_name == 'push'
    steps:
      - uses: appleboy/ssh-action@v1.0.3
        with:
          host: ${{ secrets.EC2_HOST }}
          username: ${{ secrets.EC2_USER }}
          key: ${{ secrets.EC2_SSH_KEY }}
          script: |
            cd ~/sga-sistema-distribuido
            git pull origin main
            docker compose up -d --build
\end{lstlisting}

% ─── 4. RESOLUCIÓN DE EJERCICIOS ─────────────────────────────────────────────
\section{Resolución de Ejercicios del Capítulo}

\subsection{Ejercicio 17.1: Clasificación de Pruebas en la Pirámide}
\begin{table}[H]
    \centering
    \caption{Clasificación de Pruebas en la Pirámide de Microservicios}
    \label{tab:ejercicio1}
    \begin{tabular}{p{6.5cm} p{5cm} l}
        \toprule
        \textbf{Prueba a Evaluar} & \textbf{Justificación} & \textbf{Tipo de Prueba} \\
        \midrule
        Verificar que \texttt{service.crear(dto)} lanza \texttt{ConflictoException} & Lógica de negocio aislada en memoria sin red ni base de datos. & \textbf{Unitaria (Mockito)} \\[0.2cm]
        Crear estudiante vía \texttt{POST /api} y verificar en PostgreSQL & Prueba el contexto Spring y persistencia real en Docker. & \textbf{Integración (Testcontainers)} \\[0.2cm]
        Simular 1000 usuarios en semana de exámenes & Evalúa concurrencia y saturación de recursos. & \textbf{Performance (Locust)} \\[0.2cm]
        Verificar que \texttt{svc-notas} cumple contrato de \texttt{svc-estudiantes} & Valida compatibilidad entre productor y consumidor. & \textbf{Contrato (Pact)} \\[0.2cm]
        Flujo completo: login $\rightarrow$ buscar $\rightarrow$ registrar nota & Simula interacción completa de usuario por Gateway. & \textbf{E2E (Playwright)} \\
        \bottomrule
    \end{tabular}
\end{table}

\subsection{Ejercicio 17.2: Diagnóstico de Métricas en Prometheus}
\textbf{Enunciado:} En Grafana se observa que \texttt{http\_server\_requests\_seconds\_p99} sube de 200ms a 3000ms a las 18:00 todos los días. ¿Qué podría estar ocurriendo?\\
\textbf{Diagnóstico Técnico:}
\begin{enumerate}
    \item \textbf{Pico de concurrencia al cierre de jornada laboral:} Múltiples docentes registrando calificaciones simultáneamente.
    \item \textbf{Tarea programada (Batch Job):} Proceso cron de consolidación o respaldo de base de datos a las 18:00 que satura el pool de conexiones HikariCP.
    \item \textbf{Pausa de Garbage Collection (JVM):} Ciclo de recolección mayor de memoria que suspende hilos de ejecución transaccionales.
    \item \textbf{Expiración de caché Redis:} Invalidación simultánea de claves que forza consultas directas a la base de datos PostgreSQL.
\end{enumerate}
\textbf{Plan de Acción:} Correlacionar en Grafana las métricas de conexiones activas a base de datos, consumo de CPU y registros de error en ese intervalo.

% ─── 5. CONCLUSIONES ─────────────────────────────────────────────────────────
\section{Conclusiones}

El diseño y desarrollo de la arquitectura distribuida para el módulo central \texttt{sga-principal} permitió validar la eficacia de desacoplar la lógica de dominio respecto a la infraestructura mediante Clean Architecture y principios SOLID. La integración de la suite de pruebas unitarias con JUnit 5 y Mockito demostró que es posible certificar la integridad de los casos de uso en menos de 2 segundos, reduciendo drásticamente el tiempo de retroalimentación en el ciclo de desarrollo. Asimismo, la orquestación del pipeline CI/CD con GitHub Actions y la instrumentación de observabilidad mediante Spring Boot Actuator y Prometheus garantizan que el sistema no solo sea funcional, sino altamente disponible, medible y resiliente ante fallos en producción. La experiencia adquirida en la gestión de bases de datos y la comunicación gRPC de alto rendimiento sienta las bases profesionales para liderar proyectos distribuidos de escala empresarial.

\end{document}
